% Extracción quirúrgica del propietario V2 05e: líneas 68--110.
% El resto permanece en este archivo como procedencia, pero no se publica.
\iffalse
\chapter[Macroestructura determinantal interna]{Macroestructura determinantal interna generada por el continuo discreto}
\label{chap:detint-macroestructura}

\begin{thesisbox}
La rama \({\rm det}\) parte exclusivamente de historias supervivientes de
\(\mathcal C_{\rm cont}^{\rm disc}\). El nervio de extensiones, la diagonal
Alexander--Whitney, el complejo local, la unidad, las dos publicaciones, el
filler, la reducción terminal y la torre \(3\)-ádica se construyen antes de
cualquier evaluación posterior. Las condiciones de balance y aciclicidad se
incorporan al dominio tipado y nunca se presuponen silenciosamente.
\end{thesisbox}

\section{Dominio determinantal superviviente}

Sea
\begin{equation}\label{eq:detint-an}
 A_n:=\mathbf Z/3^n\mathbf Z,
 \qquad
 \rho_{n+1,n}:A_{n+1}\longrightarrow A_n.
\end{equation}
Para \(x\in\mathcal C_{\rm cont}^{\rm disc}\), sea
\(\mathsf T_m(x)\) la categoría finita de prefijos TPK de profundidad \(m\)
compatibles con \(x\). Cada objeto conserva hoja, orientación, residuo,
cociente, carry entero, memoria, frontera, ruta, multisección, supervivencia
y terminal; sus morfismos son las extensiones admisibles que conservan esas
coordenadas.

De cada objeto \(y\in\mathsf T_m(x)\) se extraen los enteros internos
\begin{equation}\label{eq:detint-c-v-lifts}
 \widetilde c_y,\widetilde v_y\in\mathbf Z,
 \qquad
 c_{y,n}:=\widetilde c_y\bmod3^n,
 \qquad
 v_{y,n}:=\widetilde v_y\bmod3^n.
\end{equation}
La primera coordenada es el carry no orientado y la segunda, el coeficiente
orientado de frontera. La reversión satisface
\begin{equation}\label{eq:detint-or-cv}
 c_{\bar y,n}=c_{y,n},
 \qquad
 v_{\bar y,n}=-v_{y,n}.
\end{equation}

Se escriben
\begin{equation}\label{eq:detint-z3-k3}
 \mathbf Z_3:=\varprojlim_n A_n,
 \qquad
 K_3:=\operatorname{Frac}(\mathbf Z_3).
\end{equation}
La restricción determinantal se define por
\begin{equation}\label{eq:detint-dominio}
\begin{split}
 \mathcal C_{\rm cont}^{\rm det}:=
 \bigl\{x\in\mathcal C_{\rm cont}^{\rm disc}:{}&
 x\text{ sobrevive a toda profundidad};\\
 &\widehat P_x^{\rm red}\text{ es finito y basado};\\
 &\chi(\widehat P_x^{\rm red})=0;\\
 &H_*(\widehat P_x^{\rm red}\otimes_{\mathbf Z_3}K_3)=0;\\
 &\text{la reducción canónica produce un bloque cuadrado}\\
 &\text{estable bajo refinamiento}\bigr\}.
\end{split}
\end{equation}
Los objetos \(\widehat P_x^{\rm red}\) y el bloque cuadrado se construirán
en \cref{sec:detint-terminal}. Por tanto
\eqref{eq:detint-dominio} es una condición de pertenencia comprobable y no
una conclusión insertada en la definición de los mapas.

\fi
\section{Nervio de rutas, diagonal de Alexander--Whitney y counidad}

Se define el complejo normalizado
\begin{equation}\label{eq:detint-gmn}
 G_{m,n}(x):=
 N_*^{\rm norm}\!\left(N\mathsf T_m(x);A_n\right).
\end{equation}
Para un símplice no degenerado
\(\sigma=[y_0\to y_1\to\cdots\to y_q]\), la diagonal es
\begin{equation}\label{eq:detint-aw}
 \Delta_{\rm AW}(\sigma)
 =\sum_{j=0}^q
 [y_0\to\cdots\to y_j]\otimes
 [y_j\to\cdots\to y_q].
\end{equation}
La counidad viene dada por
\begin{equation}\label{eq:detint-counidad-g}
 \epsilon_G([y])=1,
 \qquad
 \epsilon_G(\sigma)=0\quad(q>0).
\end{equation}
En particular,
\begin{equation}\label{eq:detint-vertice-grouplike}
 \Delta_{\rm AW}[y]=[y]\otimes[y],
 \qquad
 \epsilon_G([y])=1.
\end{equation}

Si \(m'\ge m\) y \(n'\ge n\), truncación y reducción de coeficientes
inducen
\begin{equation}\label{eq:detint-rmn}
 R_{m',m}^{n',n}:G_{m',n'}(x)\longrightarrow G_{m,n}(x),
\end{equation}
y se verifican
\begin{equation}\label{eq:detint-aw-natural}
 R\partial=\partial R,
 \qquad
 (R\otimes R)\Delta_{\rm AW}=\Delta_{\rm AW}R,
 \qquad
 \epsilon_GR=\rho_{n',n}\epsilon_G.
\end{equation}
La ruta no se reduce a sus extremos: el símplice entero permanece en
\(G_{m,n}(x)\).

\iffalse
\section{Complejo local interno}

Para \(c\in A_n\), se define
\begin{equation}\label{eq:detint-p0-grados}
 P^0_{n,1}(c)=A_nf,
 \qquad
 P^0_{n,0}(c)=A_nr\oplus A_ne\oplus A_nb,
 \qquad
 P^0_{n,-1}(c)=A_ng,
\end{equation}
con diferencial
\begin{equation}\label{eq:detint-p0-diferencial}
 \partial f=3ce+b,
 \qquad
 \partial r=0,
 \qquad
 \partial e=g,
 \qquad
 \partial b=-3cg.
\end{equation}
La identidad de cadena es literal:
\begin{equation}\label{eq:detint-d2}
 \partial^2f=3c\,\partial e+\partial b=3cg-3cg=0,
 \qquad
 \partial^2e=\partial^2b=\partial^2r=0.
\end{equation}

La aumentación
\begin{equation}\label{eq:detint-epsilon-p}
 \epsilon_P:P_n^0(c)\longrightarrow A_n[0]
\end{equation}
se fija por
\[
 \epsilon_P(r)=1,
 \qquad
 \epsilon_P(e)=\epsilon_P(b)=\epsilon_P(f)=\epsilon_P(g)=0.
\]
El complejo posee la coalgebra diferencial
\begin{equation}\label{eq:detint-coalgebra-p}
 \Delta_P(r)=r\otimes r,
 \qquad
 \Delta_P(q)=q\otimes r+r\otimes q
 \quad(q\in\{e,b,f,g\}).
\end{equation}
Así, \(r\) es group-like y los otros cuatro generadores son primitivos
relativos a \(r\). Las ecuaciones
\eqref{eq:detint-p0-diferencial} prueban que \(\Delta_P\) conmuta con el
diferencial tensorial.

\section{Sistema local de extensiones y ledger de memoria}

Para una extensión \(\alpha:y\to y'\), definimos
\begin{equation}\label{eq:detint-psi}
 \Psi_\alpha:P_n^0(c_y)\longrightarrow P_n^0(c_{y'})
\end{equation}
mediante
\begin{equation}\label{eq:detint-psi-generadores}
\begin{gathered}
 \Psi_\alpha(r_y)=r_{y'},\quad
 \Psi_\alpha(e_y)=e_{y'},\quad
 \Psi_\alpha(g_y)=g_{y'},\quad
 \Psi_\alpha(f_y)=f_{y'},\\
 \Psi_\alpha(b_y)=b_{y'}+3(c_{y'}-c_y)e_{y'}.
\end{gathered}
\end{equation}
Es un mapa de complejos porque
\begin{equation}\label{eq:detint-psi-b-chain}
 \partial\Psi_\alpha(b_y)
 =-3c_{y'}g+3(c_{y'}-c_y)g
 =-3c_yg
 =\Psi_\alpha(\partial b_y),
\end{equation}
y
\begin{equation}\label{eq:detint-psi-f-chain}
 \Psi_\alpha(\partial f_y)
 =3c_ye+b+3(c_{y'}-c_y)e
 =3c_{y'}e+b
 =\partial\Psi_\alpha(f_y).
\end{equation}
Además,
\begin{equation}\label{eq:detint-psi-functorial}
 \Psi_\beta\Psi_\alpha=\Psi_{\beta\alpha}.
\end{equation}

El cambio de la coordenada orientada no se borra. Se registra por
\begin{equation}\label{eq:detint-k-alpha}
 K_\alpha:=(v_{y'}-v_y)f_{y'}.
\end{equation}
Para extensiones componibles,
\begin{equation}\label{eq:detint-k-coherencia}
 K_{\beta\alpha}=K_\beta+\Psi_\beta K_\alpha.
\end{equation}
Esta igualdad es el ledger exacto del incremento de frontera y hace visible
la memoria que una composición puramente nominal perdería.

El diagrama local completo es la construcción de Grothendieck
\begin{equation}\label{eq:detint-grothendieck-local}
 \int_{y\in\mathsf T_m(x)}P_n^0(c_y).
\end{equation}
En ella una flecha sobre \(\alpha:y\to y'\) lleva como componente lineal
\(\Psi_\alpha\) y como coherencia \(K_\alpha\). Cuando se aplica a un
elemento soportado en un vértice, se usa la notación
\begin{equation}\label{eq:detint-psi-soporte-vertice}
 \Psi_\alpha(p,[y]):=(\Psi_\alpha p,[y']).
\end{equation}
La arista \([y\to y']\) permanece en la coordenada de ruta del nervio. Esta
convención será la que se emplee en las fórmulas de naturalidad de
\(z_\pm\) y \(h_*\); no se afirma un mapa global que borre los restantes
símplices de \(G_{m,n}(x)\).

\section{Pullback, unidad, raíz, publicaciones y filler}

El pullback de complejos es
\begin{equation}\label{eq:detint-pgen}
 P^{\rm gen}_{y;m,n}
 :=P_n^0(c_y)\times_{A_n[0]}G_{m,n}(x)
 =\{(p,\gamma):\epsilon_P(p)=\epsilon_G(\gamma)\}.
\end{equation}
Su diferencial es
\begin{equation}\label{eq:detint-pgen-d}
 \partial(p,\gamma)=(\partial p,\partial\gamma).
\end{equation}
La unidad común asociada al vértice \(y\) es
\begin{equation}\label{eq:detint-unidad}
 \mathbf1_y:=(r,[y]).
\end{equation}
Sus dos proyecciones son group-like por
\eqref{eq:detint-vertice-grouplike} y
\eqref{eq:detint-coalgebra-p}; ése es el sentido preciso de unidad
group-like compatible del pullback.

La proyección raíz es el mapa de complejos
\begin{equation}\label{eq:detint-proyeccion-root}
 \varpi_{\rm root}:P^{\rm gen}_{y;m,n}\longrightarrow
 \operatorname{Root}_{y;m,n},
 \qquad
 \varpi_{\rm root}(p,\gamma):=(\epsilon_P(p)r,\gamma).
\end{equation}
Se definen las dos publicaciones internas
\begin{equation}\label{eq:detint-zpm}
 z_-(y):=(r,[y]),
 \qquad
 z_+(y):=(r+3c_yv_ye+v_yb,[y]).
\end{equation}
Ambas son ciclos:
\begin{equation}\label{eq:detint-zpm-ciclos}
 \partial z_-(y)=0,
 \qquad
 \partial z_+(y)=3c_yv_yg-3c_yv_yg=0.
\end{equation}
Sus raíces coinciden:
\begin{equation}\label{eq:detint-raiz-comun}
 \varpi_{\rm root}z_-(y)
 =(r,[y])
 =\varpi_{\rm root}z_+(y).
\end{equation}

El filler se construye directamente:
\begin{equation}\label{eq:detint-hstar}
 h_*(y):=(-v_yf,0)\in P^{\rm gen}_{y;m,n,1}.
\end{equation}
Entonces
\begin{equation}\label{eq:detint-filler-identidad}
 \partial h_*(y)
 =(-3c_yv_ye-v_yb,0)
 =z_-(y)-z_+(y).
\end{equation}
En particular,
\begin{equation}\label{eq:detint-clases-iguales}
 [z_-(y)]=[z_+(y)]
 \quad\text{en}\quad H_0(P^{\rm gen}_{y;m,n}),
\end{equation}
y la igualdad conserva su cadena testigo.

Bajo una extensión \(\alpha:y\to y'\), se cumplen
\begin{equation}\label{eq:detint-z-natural}
 z_-(y')=\Psi_\alpha z_-(y),
 \qquad
 z_+(y')-\Psi_\alpha z_+(y)=\partial K_\alpha,
\end{equation}
y
\begin{equation}\label{eq:detint-h-natural}
 h_*(y')-\Psi_\alpha h_*(y)=-K_\alpha.
\end{equation}
Por tanto, la naturalidad no identifica valores distintos de \(v\): registra
su diferencia mediante la homotopía coherente \(K_\alpha\).

\section{Orientación interna}

Sobre \(P_n^0(c)\) se define
\begin{equation}\label{eq:detint-omega-p}
 \Omega(r)=r,
 \qquad
 \Omega(q)=-q\quad(q\in\{e,b,f,g\}).
\end{equation}
Es un automorfismo de complejos y de la coalgebra relativa. En el nervio,
la inversión orientada es
\begin{equation}\label{eq:detint-omega-g}
 \mathfrak o_G[y_0\to\cdots\to y_q]
 :=(-1)^{q(q+1)/2}
 [\bar y_q\to\cdots\to\bar y_0].
\end{equation}
Junto con \eqref{eq:detint-or-cv}, estas acciones dan
\begin{equation}\label{eq:detint-or-z-h}
 \Omega z_-(y)=z_-(\bar y),
 \qquad
 \Omega z_+(c_y,v_y)=z_+(c_{\bar y},v_{\bar y}),
 \qquad
 \Omega h_*(y)=h_*(\bar y).
\end{equation}
La orientación actúa sobre ruta, complejo, publicaciones, filler y línea
determinantal; no es una etiqueta sin acción.

\section{Reducción terminal y condición de supervivencia}
\label{sec:detint-terminal}

Se elimina únicamente la raíz común:
\begin{equation}\label{eq:detint-p-reducido}
 \overline P_{y;m,n}
 :=P^{\rm gen}_{y;m,n}/A_n\mathbf1_y,
 \qquad
 \widehat P_{y;m}^{\rm red}:=\varprojlim_n\overline P_{y;m,n}.
\end{equation}
La base se ordena primero por el orden TPK de rutas y después por
\begin{equation}\label{eq:detint-orden-base}
 f\prec e\prec b\prec g.
\end{equation}
Se ejecuta el siguiente algoritmo interno.
\begin{enumerate}
 \item Buscar los coeficientes que son unidades de \(\mathbf Z_3\).
 \item Elegir el primero según \eqref{eq:detint-orden-base} y el orden de
 rutas.
 \item Cancelar el par correspondiente y registrar las operaciones
 elementales, su signo y su contribución a la orientación.
 \item Repetir hasta que no quede un pivote unitario cancelable.
\end{enumerate}
La condición de supervivencia de \eqref{eq:detint-dominio} exige que,
cofinalmente en \(m\), el resultado sea un bloque de dos términos
\begin{equation}\label{eq:detint-sx}
 \mathbf Z_3^{r_x}\xrightarrow{S_x}\mathbf Z_3^{r_x},
 \qquad
 \det S_x\ne0\text{ en }K_3,
\end{equation}
y que un refinamiento sólo añada bloques contractibles unitarios y cambios
de base registrados. Si esas condiciones fallan, la historia no pertenece a
\(\mathcal C_{\rm cont}^{\rm det}\); no se fuerza un terminal inexistente.

La ruta identidad aporta un testigo elemental. Después de retirar la raíz,
el complejo local es
\begin{equation}\label{eq:detint-testigo-contractil}
 A_nf\xrightarrow{\binom{3c}{1}}
 A_ne\oplus A_nb\xrightarrow{(1,-3c)}A_ng.
\end{equation}
Es contractible mediante
\begin{equation}\label{eq:detint-contraccion-local}
 s(g)=e,
 \qquad
 s(ae+\beta b)=\beta f,
 \qquad
 s(f)=0.
\end{equation}
Así, las condiciones terminales admiten al menos el testigo identidad
interno y no describen un subdominio formalmente vacío.

\section{Smith, orden \texorpdfstring{\(3\)}{3}-ádico y unidad inicial}

En bases orientadas, una diagonalización por valoración tiene la forma
\begin{equation}\label{eq:detint-smith}
 U_xS_xV_x
 =\operatorname{diag}(3^{a_1}u_1,\ldots,3^{a_r}u_r),
 \qquad
 a_i\in\mathbf N,quad u_i\in\mathbf Z_3^\times.
\end{equation}
Se define
\begin{equation}\label{eq:detint-orden-d}
 d_x:=v_3(\det S_x)=\sum_{i=1}^r a_i
\end{equation}
y la unidad inicial
\begin{equation}\label{eq:detint-leading-unit}
 \lambda_x:=3^{-d_x}\det S_x\in\mathbf Z_3^\times.
\end{equation}
La coordenada de orientación trivializa la línea determinantal. Si una
transición produce
\begin{equation}\label{eq:detint-estabilidad-bloques}
 S_x\oplus I_q=U S_{x'}V,
\end{equation}
la orientación se transporta por el factor inverso de \(\det U\det V\).
La unidad orientada \(\lambda_x^{\rm or}\) es entonces natural y no cambia
al añadir bloques contractibles.

En precisión finita,
\begin{equation}\label{eq:detint-smith-finito}
 S_{x,n}=S_x\bmod3^n,
 \qquad
 a_i^{(n)}=\min(a_i,n).
\end{equation}
El terminal es, por tanto, la familia compatible de todas las precisiones y
no un entero aislado.

\section{Torre de Bockstein y acarreo de orden superior}

Sea
\[
 C_x^1\xrightarrow{S_x}C_x^0
\]
el bloque terminal. Si una clase \([\bar y]\) admite un lift
\(y\in C_x^1\) con \(S_xy\in3^qC_x^0\), definimos
\begin{equation}\label{eq:detint-beta-q}
 \beta_q[\bar y]
 :=\left[3^{-q}S_xy\bmod3\right].
\end{equation}
El cambio de lift altera el representante por una frontera de la página
correspondiente, de modo que \(\beta_q\) está bien definido. El carry
superior es
\begin{equation}\label{eq:detint-carry-superior}
 \operatorname{Carr}_q(y):=3^{-q}S_xy.
\end{equation}
No sustituye el carry TPK de \eqref{eq:detint-c-v-lifts}; registra su
prolongación en la torre de coeficientes.

Para un bloque \(3^{a_i}u_i\), la clase persiste hasta la página \(a_i\),
y en esa página el diferencial no nulo es multiplicación por
\(\bar u_i\in\mathbf F_3^\times\). Por consiguiente,
\begin{equation}\label{eq:detint-bock-longitudes}
 \{\text{longitudes de supervivencia}\}=\{a_1,\ldots,a_r\},
 \qquad
 d_x=\sum_i a_i.
\end{equation}
El producto orientado de las trivializaciones de página satisface
\begin{equation}\label{eq:detint-bock-leading}
 \Theta_{\rm Bock}(x)=\overline{\lambda_x^{\rm or}}
 \in\mathbf F_3^\times.
\end{equation}
La prueba se hace en cada bloque diagonal y se transporta por los cambios
unimodulares registrados. Smith, Bockstein y unidad inicial son así tres
lecturas del mismo terminal.

\section{Las trece coordenadas de la restricción determinantal}

Para cada \((m,n)\), se define
\begin{equation}\label{eq:detint-d-trece}
 D_{{\rm det},m,n}(x):=
 (\mathcal F,\operatorname{Ext}_\Gamma,\operatorname{Rel},
 \mathcal N,\operatorname{or},\operatorname{Carr},\operatorname{Mem},
 \partial,\gamma,\operatorname{Msect},\operatorname{Surv},
 \operatorname{Term},\mathcal L)_{{\rm det},m,n}(x),
\end{equation}
con el contenido siguiente.
\begin{enumerate}
 \item La fibra es
 \[
  \mathcal F=(A_n,\mathsf T_m(x),
  \{P_n^0(c_y)\}_y,\{P^{\rm gen}_{y;m,n}\}_y).
 \]
 \item \(\operatorname{Ext}_\Gamma\) contiene todos los morfismos
 \(\alpha\), los mapas \(\Psi_\alpha\), las homotopías \(K_\alpha\) y
 los mapas \(R_{m',m}^{n',n}\).
 \item \(\operatorname{Rel}\) contiene \(\partial^2=0\), la ecuación de
 pullback, AW, counidad, \(\Psi_\beta\Psi_\alpha=\Psi_{\beta\alpha}\) y
 \(K_{\beta\alpha}=K_\beta+\Psi_\beta K_\alpha\).
 \item
 \(\mathcal N=(m,n,\widetilde c_y,\widetilde v_y,c_{y,n},v_{y,n},
 r_x,a_1,\ldots,a_{r_x},d_x)\).
 \item \(\operatorname{or}\) contiene \(y\mapsto\bar y\),
 \(\mathfrak o_G\), \(\Omega\) y la orientación de la línea
 determinantal.
 \item \(\operatorname{Carr}\) conserva separadamente el carry entero,
 su reducción \(c_{y,n}\) y todos los carries superiores
 \(\operatorname{Carr}_q\).
 \item \(\operatorname{Mem}\) contiene
 \[
  y_0\to\cdots\to y_m,
  \quad\{(\widetilde c_{y_j},\widetilde v_{y_j},\mathbf1_{y_j})\}_{j=0}^m,
  \quad\{K_{\alpha_j}\}.
 \]
 \item
 \(\partial=(\partial_G,\partial_P,\partial_{P^{\rm gen}},
 \partial_{\overline P},h_*)\).
 \item \(\gamma=[y_0\to\cdots\to y_q]\) conserva el símplice orientado
 completo.
 \item \(\operatorname{Msect}\) contiene todas las extensiones
 compatibles, todos los lifts \(3\)-ádicos y todas las presentaciones
 unimodulares del mismo terminal orientado; no elige una a posteriori.
 \item \(\operatorname{Surv}\) registra profundidad arbitraria,
 compatibilidad de coeficientes, balance, aciclicidad sobre \(K_3\),
 estabilidad por bloques unitarios y permanencia de la raíz.
 \item
 \[
  \operatorname{Term}=(\widehat P_x^{\rm red},S_x,
  \{a_i,u_i\},d_x,\lambda_x^{\rm or},
  \{\beta_q\},\Theta_{\rm Bock}).
 \]
 \item El lector interno es
 \[
  \mathcal L_{\rm det}^{\rm int}
  =(\mathbf1_y,
  \varpi_{\rm root}z_-=\varpi_{\rm root}z_+,
  [z_-]=[z_+],d_x,\lambda_x^{\rm or},\Theta_{\rm Bock}).
 \]
\end{enumerate}

Las transiciones satisfacen, coordenada por coordenada,
\begin{equation}\label{eq:detint-naturalidad-trece}
 u^{\rm det}_{(m',n'),(m,n)}D_{{\rm det},m',n'}(x')
 =D_{{\rm det},m,n}(\tau_{m',m}x').
\end{equation}

\section{Ensamblador, publicador y cadena no ablativa}

Se define
\begin{equation}\label{eq:detint-graph-d}
 \mathcal G_{\rm det}:=\operatorname{Graph}(D_{\rm det}),
 \qquad
 \operatorname{Res}_{\rm det}(x):=(x,D_{\rm det}x).
\end{equation}
Sea \(\chi_{\rm det}(x)\) la multisección completa de gérmenes
\[
 (y,\widetilde c_y,\widetilde v_y,\operatorname{or}_y,\gamma_y)
\]
compatibles a toda profundidad. El objeto dependiente es
\begin{equation}\label{eq:detint-x-dependiente}
 X_{\chi,{\rm det}}
 :=\{(\chi_{\rm det}(x),x,D_{\rm det}x):
 x\in\mathcal C_{\rm cont}^{\rm det}\},
\end{equation}
y
\begin{equation}\label{eq:detint-prx}
 \operatorname{pr}_{X,{\rm det}}(x,D_{\rm det}x)
 :=(\chi_{\rm det}(x),x,D_{\rm det}x).
\end{equation}

El ensamblador crudo
\begin{equation}\label{eq:detint-ensamblador-crudo}
 a_{\rm det}:X_{\chi,{\rm det}}\longrightarrow
 \mathscr M_{\rm det}^{\rm amb}
\end{equation}
produce
\begin{equation}\label{eq:detint-ensamblador-salida}
\begin{split}
 a_{\rm det}(z)=\bigl(&
 \{G_{m,n},\Delta_{\rm AW},\epsilon_G\}_{m,n};\\
 &\{P_n^0(c),\Psi_\alpha,K_\alpha,P^{\rm gen},
 \mathbf1,z_-,z_+,h_*\}_{m,n};\\
 &\{\widehat P^{\rm red},S,d,\lambda^{\rm or},
 \beta_q,\Theta_{\rm Bock}\}\bigr).
\end{split}
\end{equation}
El macroobjeto conserva su entrada:
\begin{equation}\label{eq:detint-graph-a}
 \mathfrak M_{\rm det}:=\operatorname{Graph}(a_{\rm det}),
 \qquad
 \mathcal A_{\rm det}(z):=(z,a_{\rm det}z).
\end{equation}

El publicador crudo
\begin{equation}\label{eq:detint-publicador-crudo}
 p_{\rm det}:\mathfrak M_{\rm det}\longrightarrow\mathscr Y_{\rm det}
\end{equation}
publica conjuntamente el certificado
\begin{equation}\label{eq:detint-certificado-publicado}
\begin{gathered}
 \partial^2=0,quad \text{AW y counidad},quad
 \mathbf1_y\text{ group-like compatible},\\
 \partial z_\pm=0,quad
 \varpi_{\rm root}z_-=\varpi_{\rm root}z_+,quad
 \partial h_*=z_--z_+,\\
 \Psi_\beta\Psi_\alpha=\Psi_{\beta\alpha},quad
 K_{\beta\alpha}=K_\beta+\Psi_\beta K_\alpha,\\
 d=v_3(\det S)=\sum_i a_i,quad
 \Theta_{\rm Bock}=\overline{\lambda^{\rm or}},quad
 \text{naturalidad completa bajo refinamiento}.
\end{gathered}
\end{equation}
Finalmente,
\begin{equation}\label{eq:detint-graph-p}
 \mathcal C_{\rm det}^{\rm HMT}:=\operatorname{Graph}(p_{\rm det}),
 \qquad
 \mathcal P_{\rm det}(M):=(M,p_{\rm det}M).
\end{equation}
La cadena íntegra es
\begin{equation}\label{eq:detint-cadena-completa}
 \mathcal C_{\rm cont}^{\rm disc}\supset
 \mathcal C_{\rm cont}^{\rm det}
 \xrightarrow{\operatorname{Res}_{\rm det}}\mathcal G_{\rm det}
 \xrightarrow{\operatorname{pr}_{X,{\rm det}}}X_{\chi,{\rm det}}
 \xrightarrow{\mathcal A_{\rm det}}\mathfrak M_{\rm det}
 \xrightarrow{\mathcal P_{\rm det}}\mathcal C_{\rm det}^{\rm HMT}.
\end{equation}
Cada flecha tiene una inversa sobre su imagen dada por una proyección de
coordenadas. No se sustituye historia por valor, multisección por selector,
complejo por determinante ni prueba por nombre.

\section{Teorema de realización determinantal interna}

\begin{theorem}[Realización determinantal interna completa]
\label{thm:detint-realizacion}
Para todo \(x\in\mathcal C_{\rm cont}^{\rm det}\), la cadena
\eqref{eq:detint-cadena-completa} construye un macroobjeto natural en
profundidad y precisión. Contiene dos ciclos con raíz común, un filler
explícito que los identifica, una unidad group-like compatible y un terminal
\(3\)-ádico cuyo orden, forma de Smith y torre de Bockstein producen el mismo
lector orientado. Las trece coordenadas de \eqref{eq:detint-d-trece}
permanecen recuperables en la salida.
\end{theorem}

\begin{proof}
Las ecuaciones \eqref{eq:detint-p0-diferencial}--
\eqref{eq:detint-d2} prueban que \(P_n^0(c)\) es un complejo. Las dos
aumentaciones son mapas de complejos, por lo que el pullback está tipado.
Alexander--Whitney es natural y hace group-like a cada vértice; junto con
\(r\), esto da la unidad común. El cálculo directo
\eqref{eq:detint-zpm-ciclos}--\eqref{eq:detint-filler-identidad} prueba que
las publicaciones son ciclos, tienen la misma raíz y están unidas por
\(h_*\). Las fórmulas de \(\Psi_\alpha\) prueban functorialidad, y
\(K_\alpha\) prueba la naturalidad homotópica conservando el incremento de
memoria. La definición de \(\mathcal C_{\rm cont}^{\rm det}\) garantiza que
la reducción terminal produce un bloque cuadrado sin postularlo fuera del
dominio. Smith sobre \(\mathbf Z_3\) da \(d=\sum_i a_i\); el cálculo de cada
bloque \(3^{a_i}u_i\) da
\(\Theta_{\rm Bock}=\overline{\lambda^{\rm or}}\). Truncación, reducción y
adición de bloques unitarios conservan estas identidades. Por último, las
tres construcciones por graph retienen la entrada como primera coordenada,
de modo que la composición completa es no ablativa.
\end{proof}

\section{Procedencia y falsadores}

\paragraph{Procedencia.}
El nervio TPK, Alexander--Whitney, la counidad, la unidad común, la forma
normal \(\partial f=3ce+b\), \(\partial e=g\),
\(\partial b=-3cg\), la raíz común, el filler y la arquitectura de torre
tienen estatuto \texttt{ARQUITECTURA\_AUTORAL\_PREEXISTENTE}. El cálculo
literal de los dos ciclos, su raíz y la identidad
\(\partial h_*=z_--z_+\) tiene estatuto \texttt{RESULTADO\_RECUPERADO}. El
sistema local \(\Psi_\alpha/K_\alpha\), el subdominio superviviente
balanceado, el terminal exclusivamente \(3\)-ádico y el empaquetado por
graphs tienen estatuto \texttt{FORMALIZACION\_NUEVA}; formalizan sin cambiar
la procedencia conceptual del aparato.

\begin{ablation}[Falsadores internos de la rama determinantal]
\label{abl:detint-falsadores}
La construcción queda falsada en el nivel correspondiente si ocurre alguno
de los hechos siguientes.
\begin{enumerate}
 \item Falla \(\partial^2=0\); entonces no existe el complejo local.
 \item Alexander--Whitney no conmuta con refinamiento, la counidad falla o
 el vértice deja de ser group-like.
 \item \(z_\pm\) no son ciclos, sus raíces difieren o falla
 \(\partial h_*=z_--z_+\).
 \item \(\Psi_\alpha\) no es mapa de complejos o falla
 \(K_{\beta\alpha}=K_\beta+\Psi_\beta K_\alpha\); en ese caso se ha
 perdido carry o memoria.
 \item El complejo reducido no es balanceado o no es acíclico sobre
 \(K_3\); la historia queda fuera del dominio terminal tipado.
 \item Las longitudes de Bockstein no coinciden con los \(a_i\), o
 \(\Theta_{\rm Bock}\ne\overline{\lambda^{\rm or}}\).
 \item El refinamiento modifica \(d\) o \(\lambda^{\rm or}\) fuera de
 bloques unitarios y cambios de orientación registrados.
 \item Se omite una de las trece coordenadas: \emph{objeto sustituido;
 conclusión no transferible}.
\end{enumerate}
Estos falsadores comprueban el objeto interno construido; no lo reemplazan
por un control posterior.
\end{ablation}
\fi
